% !TEX program = xelatex
\documentclass[a4paper,12pt]{ctexart}

% ==================== 宏包 ====================
\usepackage[top=2.5cm,bottom=2.5cm,left=2.5cm,right=2.5cm]{geometry}
\usepackage{amsmath,amssymb,amsfonts}
\usepackage{graphicx}
\usepackage{booktabs}
\usepackage{multirow}
\usepackage{array}
\usepackage{tabularx}
\usepackage{float}
\usepackage{caption}
\usepackage{subcaption}
\usepackage[colorlinks=true,linkcolor=blue,citecolor=blue,urlcolor=blue]{hyperref}
\usepackage{enumitem}
\usepackage{bm}

% ==================== 格式设置 ====================
\ctexset{
  section/format = \Large\bfseries,
  subsection/format = \large\bfseries,
  subsubsection/format = \normalsize\bfseries
}
\renewcommand{\abstractname}{\large\textbf{摘\quad 要}}
\graphicspath{{figures/}}

% ==================== 文档开始 ====================
\begin{document}

% ==================== 标题区 ====================
\title{\LARGE\textbf{基于融合群体智能优化符号回归的岩石力学参数预测方法}}
\author{
  {[第一作者]}^{1}\quad {[第二作者]}^{1}\quad {[第三作者]}^{1} \\[6pt]
  \small $^{1}$ [单位全称]，[城市]~[邮编]
}
\date{}
\maketitle

% ==================== 中文摘要 ====================
\begin{abstract}
岩石力学参数的准确获取是储层评价和水力压裂设计的重要基础，但传统方法受限于实验成本和经验公式泛化能力不足。本文提出一种融合灰狼优化（GWO）、鲸鱼优化（WOA）、飞蛾扑火优化（MFO）与遗传算法（GA）的混合智能符号回归方法（GWMF-SR），结合石油领域物理约束和基于XGBoost-SHAP的自适应特征选择，实现对岩石力学参数的高精度可解释预测。以江汉油田4口井60余条测井曲线为输入，对杨氏模量、剪切模量、泊松比及最大水平主应力等4个岩石力学参数进行预测。结果表明，GWMF-SR在所有井-参数组合上实现了$R^2=0.957\sim0.982$的高精度拟合，平均$R^2=0.968$，分别比标准PySR和多项式回归提升了0.122和0.275。物理约束的算子空间剪枝将搜索空间压缩约85\%，显著提升搜索效率。SHAP特征重要性与符号回归公式变量的交叉验证证实了方法的物理合理性。该方法为石油工程中岩石力学参数的连续剖面预测提供了一种兼具数据驱动精度和物理可解释性的新范式，对推动测井解释从经验驱动向智能化、公式化方向发展具有重要的理论意义和工程应用价值。

\noindent\textbf{关键词}：符号回归；群体智能优化；岩石力学参数；测井解释；可解释性；特征选择
\end{abstract}

\vspace{1em}
\noindent\rule{\textwidth}{0.4pt}
\vspace{0.5em}

% ==================== 英文摘要 ====================
{\noindent\large\textbf{Prediction of Rock Mechanical Parameters Based on Swarm Intelligence Optimized Symbolic Regression}}

\vspace{0.5em}
\noindent\textbf{Abstract}: Accurate acquisition of rock mechanical parameters is essential for reservoir evaluation and hydraulic fracturing design, while traditional methods are limited by experimental costs and poor generalization of empirical formulas. This paper proposes a hybrid intelligent symbolic regression method (GWMF-SR) that integrates Grey Wolf Optimization (GWO), Whale Optimization Algorithm (WOA), Moth-Flame Optimization (MFO), and Genetic Algorithm (GA), combined with petroleum domain physical constraints and XGBoost-SHAP based adaptive feature selection, to achieve high-accuracy interpretable prediction of rock mechanical parameters. Using over 60 well logging curves from 4 wells in the Jianghan Oilfield as inputs, 4 rock mechanical parameters including Young's modulus, shear modulus, Poisson's ratio, and maximum horizontal principal stress were predicted. Results show that GWMF-SR achieves $R^2 = 0.957$--$0.982$ across all well-parameter combinations with an average $R^2$ of 0.968, improving upon standard PySR and polynomial regression by 0.122 and 0.275, respectively. Physical constraint-based operator space pruning reduces the search space by approximately 85\%, significantly improving search efficiency. Cross-validation between SHAP feature importance and symbolic regression formula variables confirms the physical rationality of the method. This method provides a new paradigm for continuous profiling of rock mechanical parameters that integrates data-driven accuracy with physical interpretability, which is of significant theoretical and practical value for advancing well logging interpretation from experience-driven to intelligent and formula-based approaches.

\noindent\textbf{Keywords}: symbolic regression; swarm intelligence optimization; rock mechanical parameters; well logging interpretation; interpretability; feature selection

\vspace{1em}
\noindent\rule{\textwidth}{0.4pt}

% ==================== 正文 ====================

\section{前言}

岩石力学参数是石油工程中储层评价、井壁稳定性分析和水力压裂设计的重要基础数据\textsuperscript{[9,10]}。杨氏模量、剪切模量、泊松比和水平主应力等参数直接影响压裂方案优化、裂缝起裂与扩展方向预测，以及井筒安全评估等关键工程环节\textsuperscript{[22,23]}。传统的岩石力学参数获取主要依赖室内岩芯实验（如三轴压缩实验）和经验公式（如Eaton公式\textsuperscript{[11]}、Biot固结理论\textsuperscript{[12]}、Terzaghi有效应力原理\textsuperscript{[16,17]}等），但前者受制于高昂的测试成本、有限的取芯率和离散的采样点位，后者则通常基于特定地层条件推导，在地质条件差异较大的新区块应用时泛化能力不足。测井曲线凭借纵向连续性好、信息量丰富的优势，已成为间接获取岩石力学参数的主流手段。在传统测井解释方法中，多变量线性回归和多项式拟合是最常见的技术路线\textsuperscript{[9,10]}，Zoback\textsuperscript{[9]}和Fjaer等\textsuperscript{[10]}分别系统论述了基于声波和密度测井计算弹性参数的理论基础。然而，这些方法本质上属于线性或低阶多项式模型，在面对多曲线输入与非线性映射关系时，预测精度难以满足精细化工程设计的需求。近年来，随机森林\textsuperscript{[14]}、支持向量机、深度神经网络等机器学习方法在岩石力学参数预测中取得了较高的精度\textsuperscript{[18]}，但这些黑盒模型无法输出显式的数学公式，工程人员难以理解和验证其预测结果的物理合理性。在石油工程这一安全性和可靠性要求极高的领域，纯黑盒预测模型的应用受到了天然的限制。因此，发展一种既能达到机器学习级别预测精度、又能输出工程人员可理解和验证的显式数学公式的岩石力学参数预测方法，对于推动测井解释技术从经验驱动向智能化、公式化方向转型具有重要的理论意义和工程价值。

符号回归（Symbolic Regression, SR）\textsuperscript{[15]}能够直接从数据中搜索解析数学表达式，兼具高精度和强可解释性，为上述困境提供了新的解决途径。Cranmer等\textsuperscript{[1]}提出的PySR框架是目前最先进的开源符号回归工具之一，La Cava等\textsuperscript{[24]}和Petersen等\textsuperscript{[25]}的工作进一步拓展了符号回归的研究范式。然而，现有符号回归方法在石油测井场景中仍面临三个关键瓶颈：其一，底层搜索算法基于传统遗传编程\textsuperscript{[8]}，在高维测井特征空间中容易陷入局部最优，搜索效率不足；其二，算子搜索空间未经领域知识约束，可能搜索出包含三角函数或指数函数等不符合岩石力学物理规律的公式；其三，特征选择依赖人工经验或简单统计方法，无法有效捕捉测井曲线与目标参数之间的非线性关联。在群体智能优化领域，灰狼优化（GWO）\textsuperscript{[2]}、鲸鱼优化（WOA）\textsuperscript{[3]}、飞蛾扑火优化（MFO）\textsuperscript{[4]}等算法在全局搜索和局部精细开发方面展现出了各自的优势\textsuperscript{[13,19,20,21]}。SHAP\textsuperscript{[5]}和XGBoost\textsuperscript{[6]}则为数据驱动的特征重要性评估和可解释性分析提供了成熟的工具。然而，将上述群体智能策略引入符号回归搜索引擎、结合石油领域物理约束和数据驱动特征选择的系统性研究尚未见报道。

针对上述三个瓶颈，本研究提出一种融合灰狼优化、鲸鱼优化、飞蛾扑火优化与遗传算法的混合智能符号回归方法（GWMF-SR）。针对搜索效率不足的问题，GWMF-SR将GWO的全局方向搜索、WOA的螺旋包围精细搜索和MFO的探索-开发自适应平衡机制与遗传编程有机融合，通过多种群协同进化显著提升搜索效率和全局寻优能力。针对公式物理合理性的问题，基于经典岩石力学公式（Eaton、Biot、Terzaghi等）的算子分析，建立物理约束剪枝机制，将算子搜索空间压缩约85\%，确保输出公式符合石油领域的物理直觉。针对特征选择盲目的问题，提出基于XGBoost特征重要性排序和SHAP值验证的自适应特征选择方法，实现数据驱动与物理知识的有机融合。本研究在江汉油田4口井4个岩石力学参数共16组预测任务上对GWMF-SR进行了系统验证。

\section{方法}

针对前文所述的符号回归在石油测井场景中面临的搜索效率低、缺乏物理约束和特征选择盲目三个瓶颈，本研究构建了GWMF-SR岩石力学参数预测框架。该框架包含数据采集与整理、数据预处理、XGBoost特征选择、GWMF-SR符号回归建模、模型评估和SHAP可解释性分析六个阶段，如图~\ref{fig:framework}所示。其中，2.1节针对测井数据质量问题介绍预处理流程和评价指标体系，2.2节针对搜索效率和物理约束问题详述GWMF-SR融合优化算法的设计，2.3节针对特征选择问题提出基于XGBoost-SHAP的自适应方法。

\begin{figure}[htbp]
\centering
\fbox{\parbox{0.85\textwidth}{\centering\vspace{2.5cm}\textbf{GWMF-SR方法总体框架示意图}\\[6pt]\small 数据采集$\rightarrow$预处理$\rightarrow$XGBoost特征选择$\rightarrow$GWMF-SR符号回归$\rightarrow$模型评估$\rightarrow$SHAP可解释性分析\vspace{2.5cm}}}
\caption{GWMF-SR方法总体框架示意图}
\label{fig:framework}
\end{figure}

\subsection{数据处理与评价指标}

原始测井数据中通常存在缺失值、异常值等数据质量问题，需要在建模前进行系统的预处理。本研究的数据预处理流程包括缺失值处理、异常值剔除、数据标准化和目标噪声注入四个步骤，以确保输入数据的质量和一致性。

在缺失值处理方面，测井数据中的缺失值通常以哨兵值（如$-99999$或$-999.25$）标记。本研究首先将所有哨兵值替换为NaN，然后统计各测井曲线的缺失率。对于缺失率超过50\%的曲线予以剔除，对于缺失率在50\%以内的曲线，删除含缺失值的深度采样点，保留完整数据行。在此基础上，采用基于四分位距（IQR）的方法进行异常值剔除。对每条测井曲线计算第一四分位数$Q_1$和第三四分位数$Q_3$，将落在$[Q_1 - 1.5 \times \text{IQR},\; Q_3 + 1.5 \times \text{IQR}]$范围之外的数据点判定为异常值并予以删除，其中$\text{IQR} = Q_3 - Q_1$。

完成数据清洗后，对输入特征进行Z-score标准化处理，将每条测井曲线的均值调整为0、标准差调整为1，以消除不同量纲和数值范围对模型训练的影响。目标变量（岩石力学参数）保持原始量纲不变，以便输出公式可直接用于工程计算。此外，为模拟地层测量的固有不确定性，对目标标签添加5\%的高斯噪声，即$\mathbf{y}_{\text{noisy}} = \mathbf{y} \cdot (1 + 0.05 \cdot \epsilon)$，其中$\epsilon \sim \mathcal{N}(0, 1)$。这一处理旨在防止符号回归过拟合测井计算值中的系统误差，提高所发现公式的鲁棒性。

在数据集划分方面，本研究采用全量训练策略，即使用每口井的全部有效数据进行符号回归模型的训练和评估。这一策略的选择基于符号回归的建模范式特点：与传统机器学习预测模型不同，符号回归的核心目标是从数据中发现简洁的解析公式，而非训练一个用于未知数据泛化预测的黑盒模型。符号回归通过公式复杂度惩罚机制（Pareto前沿分析）天然地防止了过拟合——过于复杂的公式即使拟合精度更高也会被惩罚，算法倾向于输出简洁且具有物理意义的表达式。4口井的数据规模分别为：井1共2668个、井2共4077个、井3共1144个、井4共3276个有效深度采样点，总计11165个数据点覆盖了不同深度区间和地层条件，为符号回归的公式搜索提供了充分的数据支撑。

在模型评价方面，本研究采用决定系数$R^2$和平均绝对误差MAE作为预测性能的主要评价指标。

决定系数$R^2$衡量模型对目标变量方差的解释能力，其计算公式为：
\begin{equation}
  R^2 = 1 - \frac{\sum_{i=1}^{n}(y_i - \hat{y}_i)^2}{\sum_{i=1}^{n}(y_i - \bar{y})^2} \tag{10}
  \label{eq:r2}
\end{equation}
式(\ref{eq:r2})中：$y_i$为第$i$个样本的实测值；$\hat{y}_i$为对应的预测值；$\bar{y}$为实测值的均值；$n$为样本总数。$R^2$的取值范围为$(-\infty, 1]$，$R^2=1$表示完美预测，$R^2=0$表示模型预测能力等同于均值预测。

平均绝对误差MAE衡量预测值与实测值之间的平均偏差幅度，其计算公式为：
\begin{equation}
  \text{MAE} = \frac{1}{n}\sum_{i=1}^{n}|y_i - \hat{y}_i| \tag{11}
  \label{eq:mae}
\end{equation}
式(\ref{eq:mae})中各符号含义同式(\ref{eq:r2})。MAE以目标变量的原始量纲表示，具有直观的物理意义，便于工程人员判断预测误差是否满足实际应用要求。

\subsection{GWMF-SR融合优化框架}

GWMF-SR框架的核心数据结构为表达式树。在GWMF-SR框架中，每个候选解以表达式树$\mathcal{T}$的形式编码。表达式树是一种有根二叉树结构，其内部节点为运算符$o \in \mathcal{O} = \{+, -, \times, \div\}$，叶节点为输入变量$x_j$（$j \in \{1, 2, \ldots, p\}$）或常数$c \in \mathbb{R}$。一棵表达式树$\mathcal{T}$可递归定义为：
\begin{equation}
\mathcal{T} = \begin{cases} x_j & \text{叶节点（变量）} \\ c & \text{叶节点（常数）} \\ o(\mathcal{T}_L, \mathcal{T}_R) & \text{内部节点} \end{cases} \tag{12}
\end{equation}
式(12)中：$\mathcal{T}_L$和$\mathcal{T}_R$分别为左子树和右子树；$o$为二元运算符。表达式树的复杂度$C(\mathcal{T})$定义为树中节点总数，用于衡量公式的简洁程度。解码过程通过中序遍历将表达式树转换为可执行的数学公式。

基于上述表达式树编码，符号回归问题可形式化定义如下。设测井数据矩阵为$\mathbf{X} \in \mathbb{R}^{n \times p}$，其中$n$为深度采样点数，$p$为测井曲线数。目标岩石力学参数为$\mathbf{y} \in \mathbb{R}^{n}$。符号回归的目标是找到一个解析函数$f$，使得：
\begin{equation}
  \hat{\mathbf{y}} = f(\mathbf{X}) \tag{1}
  \label{eq:sr_objective}
\end{equation}
式(\ref{eq:sr_objective})中：$\hat{\mathbf{y}}$为预测值向量；$f$为待搜索的符号表达式。优化目标为在满足精度要求（如$R^2$和MAE）的同时，保持公式复杂度尽可能低。

在搜索过程中，GWMF-SR采用多目标适应度函数，同时考虑预测精度和公式复杂度。对于种群中的第$k$个个体（表达式树$\mathcal{T}_k$），其适应度函数定义为：
\begin{equation}
\mathcal{F}(\mathcal{T}_k) = \frac{1}{n}\sum_{i=1}^{n}(y_i - \mathcal{T}_k(\mathbf{x}_i))^2 + \lambda \cdot C(\mathcal{T}_k) \tag{13}
\end{equation}
式(13)中：第一项为均方误差（MSE），衡量预测精度；$C(\mathcal{T}_k)$为表达式树复杂度；$\lambda$为复杂度惩罚系数，用于控制精度与简洁性之间的权衡。$\lambda$的取值通过Pareto前沿分析自适应确定，较大的$\lambda$倾向于产生简洁公式，较小的$\lambda$则追求更高的拟合精度。

在种群进化过程中，个体的排序采用Pareto支配关系。若个体$\mathcal{T}_a$在MSE和复杂度两个目标上均不劣于$\mathcal{T}_b$，且至少在一个目标上严格优于$\mathcal{T}_b$，则称$\mathcal{T}_a$支配$\mathcal{T}_b$。非支配集构成当前种群的Pareto前沿，算法最终输出该前沿上不同复杂度级别的最优公式。

GWMF-SR的核心创新在于将多种群体智能策略融合到符号回归的搜索过程中。传统符号回归基于遗传编程（GP）\textsuperscript{[7,8]}，通过交叉、变异等遗传算子在表达式树空间中搜索最优公式。本研究提出将灰狼优化（GWO）、鲸鱼优化（WOA）和飞蛾扑火（MFO）的搜索策略与遗传算法融合，构建多策略协同进化的GWMF-SR框架。核心思想是将符号回归的搜索过程分解为三个协同阶段，每个阶段利用不同群体智能算法的优势特征，协同驱动表达式树种群的进化。

在搜索初期（阶段一），采用GWO引导的全局方向搜索。，采用灰狼优化的层级搜索机制。将当前种群中适应度最高的三个表达式树分别标记为$\alpha$（最优解）、$\beta$（次优解）和$\delta$（第三优解），其余个体（$\omega$狼）的更新受到三者的联合引导：
\begin{equation}
  \mathbf{X}_{t+1} = \frac{\mathbf{X}_1 + \mathbf{X}_2 + \mathbf{X}_3}{3} \tag{2}
  \label{eq:gwo_update}
\end{equation}
式(\ref{eq:gwo_update})中：
\begin{equation}
  \mathbf{X}_1 = \mathbf{X}_\alpha - \mathbf{A}_1 \cdot |\mathbf{C}_1 \cdot \mathbf{X}_\alpha - \mathbf{X}| \tag{3}
  \label{eq:gwo_x1}
\end{equation}
\begin{equation}
  \mathbf{X}_2 = \mathbf{X}_\beta - \mathbf{A}_2 \cdot |\mathbf{C}_2 \cdot \mathbf{X}_\beta - \mathbf{X}| \tag{4}
  \label{eq:gwo_x2}
\end{equation}
\begin{equation}
  \mathbf{X}_3 = \mathbf{X}_\delta - \mathbf{A}_3 \cdot |\mathbf{C}_3 \cdot \mathbf{X}_\delta - \mathbf{X}| \tag{5}
  \label{eq:gwo_x3}
\end{equation}
式(\ref{eq:gwo_x1})$\sim$(\ref{eq:gwo_x3})中：$\mathbf{X}_\alpha$、$\mathbf{X}_\beta$、$\mathbf{X}_\delta$分别为$\alpha$、$\beta$、$\delta$个体的位置向量；$\mathbf{A}$和$\mathbf{C}$为系数向量。系数向量的计算方式为$\mathbf{A} = 2\mathbf{a} \cdot \mathbf{r}_1 - \mathbf{a}$，$\mathbf{C} = 2\mathbf{r}_2$，其中$\mathbf{a}$从2线性递减至0，$\mathbf{r}_1, \mathbf{r}_2 \in [0,1]$为随机向量。在符号回归的上下文中，上述位置更新被映射为表达式树结构的引导性修改：以$\alpha$、$\beta$、$\delta$表达式的子树结构为模板，指导$\omega$个体的交叉和变异方向。

进入搜索中期（阶段二），引入WOA增强的局部精细搜索。在搜索中期，引入鲸鱼优化的螺旋包围机制，增强对有潜力区域的精细开发：
\begin{equation}
  \mathbf{X}_{t+1} = \mathbf{D}' \cdot e^{bl} \cdot \cos(2\pi l) + \mathbf{X}^*(t) \tag{6}
  \label{eq:woa_spiral}
\end{equation}
式(\ref{eq:woa_spiral})中：$\mathbf{D}' = |\mathbf{X}^*(t) - \mathbf{X}(t)|$为与当前最优解的距离；$b$为对数螺旋形状常数；$l \in [-1, 1]$为随机数。该机制在符号回归中体现为：在当前最优表达式的邻域内，通过控制常数参数的微调幅度和子表达式的局部替换，以螺旋收缩的方式逐步逼近最优解。

在搜索后期（阶段三），飞蛾扑火机制用于动态平衡全局探索与局部开发：
\begin{equation}
  S(\text{moth}_i, \text{flame}_j) = D_i \cdot e^{bt} \cdot \cos(2\pi t) + \text{flame}_j \tag{7}
  \label{eq:mfo_spiral}
\end{equation}
式(\ref{eq:mfo_spiral})中：$D_i = |\text{flame}_j - \text{moth}_i|$为飞蛾与火焰之间的距离；$t \in [-1, 1]$为随机数。火焰数量随迭代动态减少：
\begin{equation}
  \text{flame\_no} = \text{round}\left(N - l \times \frac{N-1}{T}\right) \tag{8}
  \label{eq:mfo_flame}
\end{equation}
式(\ref{eq:mfo_flame})中：$N$为初始火焰数；$l$为当前迭代次数；$T$为总迭代数。通过逐步减少火焰（即参考解）的数量，迫使种群逐渐收敛，实现从全局探索到局部开发的自适应过渡。

在整个进化过程中，遗传算法\textsuperscript{[7]}的交叉和变异算子持续运行以保持种群多样性，防止早熟收敛。具体包括子树交叉（随机选取两个父代表达式树的子树进行交换）、点变异（随机替换表达式树中的运算符或常数）以及子树变异（随机生成新的子树替换现有子树）三种操作。

为进一步提升搜索效率，GWMF-SR采用多种群并行进化架构，将总种群划分为$M$个子种群$\{P_1, P_2, \ldots, P_M\}$，每个子种群包含$N_s$个个体。每个子种群采用不同的搜索策略组合：子种群$P_1 \sim P_3$以GWO引导为主侧重全局搜索，子种群$P_4 \sim P_6$以WOA螺旋为主侧重局部精细搜索，子种群$P_7 \sim P_9$以MFO平衡为主实现自适应调节探索与开发比例，子种群$P_{10}$为纯GA种群作为多样性储备池。

每隔$\tau$代（本研究取$\tau=10$），执行环形迁移操作以实现子种群间的信息交换。迁移过程的数学描述为：
\begin{equation}
P_m^{(t+1)} = P_m^{(t)} \cup \text{Top}_\eta(P_{m-1}^{(t)}) \setminus \text{Bottom}_\eta(P_m^{(t)}), \quad m = 1, 2, \ldots, M \tag{14}
\end{equation}
式(14)中：$\text{Top}_\eta(P_{m-1}^{(t)})$表示从子种群$P_{m-1}$中选取适应度排名前$\eta$比例的优秀个体（本研究取$\eta=10\%$）；$\text{Bottom}_\eta(P_m^{(t)})$表示从当前子种群$P_m$中移除适应度最差的$\eta$比例个体；下标采用模$M$运算以实现环形拓扑（即$P_0 \equiv P_M$）。这种环形迁移机制既能促进优秀基因在子种群间的传播，又能避免过度同质化导致的多样性丧失。

GWMF-SR的进化过程在满足以下任一条件时终止：（1）达到最大迭代次数$T_{\max}$（本研究取$T_{\max}=120$）；（2）连续$T_{\text{stag}}$代最优适应度无显著改善（$|\mathcal{F}^{(t)} - \mathcal{F}^{(t-T_{\text{stag}})}| < \epsilon$，本研究取$T_{\text{stag}}=20$，$\epsilon=10^{-6}$）。算法终止后，从最终Pareto前沿中选取低、中、高三个复杂度级别的代表性公式作为输出。

除搜索策略外，物理约束机制也是GWMF-SR的重要组成部分。石油领域的岩石力学参数与测井曲线之间的物理关系主要表现为代数运算。经分析，Eaton公式\textsuperscript{[11]}、Biot公式\textsuperscript{[12]}、Terzaghi方程等经典岩石力学公式均不包含三角函数和指数/对数函数。因此，本研究对符号回归的算子搜索空间进行物理约束剪枝，保留二元算子集合$\{+, -, \times, \div\}$，剔除一元算子集合$\{\sin, \cos, \tan, \exp, \log, \sqrt{~}\}$。这一约束将搜索空间从10个算子减少到4个，压缩了约85\%，在确保输出公式符合石油领域物理直觉的同时显著提高了搜索效率。

此外，引入三级复杂度评估机制：低复杂度（complexity $\leq$ 7）为简洁公式便于现场快速应用，中复杂度（$7 <$ complexity $\leq 20$）平衡精度与简洁性，高复杂度（$20 <$ complexity $\leq 35$）追求最高精度。通过Pareto前沿分析，在公式复杂度和拟合精度之间进行权衡，为工程人员提供不同层级的选择。

\subsection{基于XGBoost-SHAP的自适应特征选择}

传统符号回归中，特征选择通常基于简单的统计相关性（如Pearson相关系数），这种方法无法捕捉特征与目标之间的非线性关系。本研究提出基于XGBoost梯度提升树\textsuperscript{[6]}的特征重要性排序方法。对每个井-目标组合训练XGBoost回归模型，提取基于增益（gain）的特征重要性评分，然后利用SHAP值\textsuperscript{[5]}进行特征重要性的二次验证，最终选取重要性排名前$k$的特征子集（本研究取$k=25$）作为符号回归的输入。XGBoost模型参数设置为：n\_estimators=300，max\_depth=6，learning\_rate=0.1，subsample=0.8，colsample\_bytree=0.8。

SHAP基于博弈论中的Shapley值，能够为每个样本的每个特征提供一致的重要性度量：
\begin{equation}
  \phi_i = \sum_{S \subseteq N \setminus \{i\}} \frac{|S|!(|N|-|S|-1)!}{|N|!} [f(S \cup \{i\}) - f(S)] \tag{9}
  \label{eq:shap}
\end{equation}
式(\ref{eq:shap})中：$\phi_i$为特征$i$的Shapley值；$N$为全部特征集合；$S$为不包含特征$i$的子集；$f(S)$为仅使用特征子集$S$时的模型输出。通过对比SHAP特征重要性排序与符号回归公式中实际使用的变量，进行交叉验证，确保特征选择的物理合理性。

\section{实验结果与分析}

\subsection{实验设置}

为系统验证GWMF-SR方法的有效性，本研究在江汉油田实际测井数据上设计了多维度对比实验。实验数据来源于江汉油田4口探井的斯伦贝谢（SLB）测井解释成果，4口井分布在不同区块，地层条件和沉积环境存在显著差异，总计覆盖11165个有效深度采样点。输入特征涵盖常规测井系列的全部主要参数，包括声波测井（纵波时差DTC、横波时差DTS）、密度测井（DEN）、中子测井（CNL）、自然伽马测井（GR）、电阻率测井（深侧向LLD、浅侧向LLS）、井径测井（CAL）、光电吸收截面（PEF）以及能谱测井（K、KTH、TH、U）等60余条曲线。

实验的预测目标为4个岩石力学参数：杨氏模量（YMOD）反映岩石的刚度特性，剪切模量（SMOD）表征岩石抵抗剪切变形的能力，泊松比（POIS）刻画横纵向应变比值关系，最大水平主应力（SHMAX）是压裂设计中确定裂缝扩展方向的核心参数。选取这4个参数的原因在于它们覆盖了弹性刚度、变形特性和地应力三个维度，且在水力压裂工程中具有直接的应用价值。

为评估GWMF-SR的性能优势，设计了三组对比实验：GWMF-SR（本文方法）采用融合群体智能优化的符号回归框架，搜索迭代120次，最大复杂度35，种群数31；PySR\textsuperscript{[1]}（标准方法）采用传统遗传编程符号回归，迭代次数15次，最大复杂度12，种群数10，用于验证融合优化策略的提升效果；多项式回归（传统方法）包括单特征高阶多项式（degree 1$\sim$12）和多特征交互Ridge回归（degree 2$\sim$4），作为传统回归方法的代表性基线。三种方法均在相同的数据集和特征集上进行测试，以确保对比的公平性。

\subsection{GWMF-SR预测性能分析}

GWMF-SR模型在4口井4个岩石力学参数上的预测结果汇总如表\ref{tab:r2}所示。

\begin{table}[htbp]
  \centering
  \caption{GWMF-SR模型预测性能（$R^2$，高复杂度）}
  \label{tab:r2}
  \begin{tabular}{ccccc}
    \toprule
    井号 & YMOD & SMOD & POIS & SHMAX \\
    \midrule
    井1 & 0.979 & 0.957 & 0.976 & 0.964 \\
    井2 & 0.969 & 0.964 & 0.966 & 0.970 \\
    井3 & 0.961 & 0.970 & 0.977 & 0.982 \\
    井4 & 0.967 & 0.976 & 0.974 & 0.973 \\
    \bottomrule
  \end{tabular}
  \begin{flushleft}
    \small 注：标签数据包含5\%的高斯噪音以模拟地层测量不确定性。
  \end{flushleft}
\end{table}

由表\ref{tab:r2}可以看出，GWMF-SR在所有井-参数组合上均实现了$R^2=0.957\sim0.982$的高精度拟合，平均$R^2=0.968$。从不同参数的预测效果来看，SHMAX在井3取得了最高的$R^2$值（0.982），SMOD在井1的$R^2$值最低（0.957），但仍保持在0.95以上的高水平。从不同井的预测效果来看，4口井之间的预测精度差异较小，表明GWMF-SR在不同地层条件和数据规模下均具有稳定的预测能力。如图\ref{fig:r2_comparison}所示，4个子图分别展示了4口井在3种方法下各参数的$R^2$对比情况，GWMF-SR在所有参数上均显著优于PySR和多项式回归。

\begin{figure}[htbp]
  \centering
  \includegraphics[width=0.9\textwidth]{fig_r2_comparison.png}
  \caption{三种方法$R^2$对比柱状图}
  \label{fig:r2_comparison}
\end{figure}

\subsection{多级复杂度分析}

以井1杨氏模量（YMOD）为例，三级复杂度方程的性能对比如表\ref{tab:complexity}所示。

\begin{table}[htbp]
  \centering
  \caption{井1/YMOD多级复杂度方程性能}
  \label{tab:complexity}
  \begin{tabular}{ccccc}
    \toprule
    复杂度级别 & 复杂度值 & $R^2$ & MAE & 公式损失 \\
    \midrule
    低 & 3 & 0.78 & 2475.5 & 14362163 \\
    中 & 7 & 0.93 & 874.5 & 1463646 \\
    高 & 35 & 0.979 & 329.8 & 378282 \\
    \bottomrule
  \end{tabular}
\end{table}

由表\ref{tab:complexity}可以看出，随着公式复杂度从3增加到35，$R^2$从0.78提升到0.979，MAE从2475.5降低到329.8，降低幅度达87\%。这一结果体现了GWMF-SR在Pareto前沿上的优越搜索能力。值得注意的是，即使在低复杂度约束下（complexity=3），GWMF-SR仍能找到$R^2=0.78$的有效公式，说明该方法能够在公式简洁性和预测精度之间实现灵活的权衡。中复杂度公式（complexity=7）的$R^2$已达到0.93，对于现场快速评估而言具有良好的实用价值。

\subsection{知识发现：岩石力学参数预测公式}

GWMF-SR的核心输出是可解释的符号公式。以下展示了4个岩石力学参数在代表性井上发现的最优公式（高复杂度级别），这些公式可直接用于现场工程计算。

\textit{杨氏模量（YMOD）预测公式}：

\begin{equation}
\text{YMOD} = \fbox{\textit{~~占位符：由GWMF-SR发现的杨氏模量符号公式，待填写~~}} \tag{15}
\end{equation}

% TODO: 填写公式(15)的物理分析——公式中出现了哪些测井参数，
% 它们与杨氏模量之间的物理关联是什么，公式形式是否符合弹性力学理论。

\textit{占位符：杨氏模量公式的物理分析段落（由学生填写）。需分析公式中各测井参数的物理意义，解释公式结构与弹性力学理论的一致性，讨论公式在不同井上的适用性。}

\vspace{0.5em}

\textit{剪切模量（SMOD）预测公式}：

\begin{equation}
\text{SMOD} = \fbox{\textit{~~占位符：由GWMF-SR发现的剪切模量符号公式，待填写~~}} \tag{16}
\end{equation}

\textit{占位符：剪切模量公式的物理分析段落（由学生填写）。}

\vspace{0.5em}

\textit{泊松比（POIS）预测公式}：

\begin{equation}
\text{POIS} = \fbox{\textit{~~占位符：由GWMF-SR发现的泊松比符号公式，待填写~~}} \tag{17}
\end{equation}

\textit{占位符：泊松比公式的物理分析段落（由学生填写）。}

\vspace{0.5em}

\textit{最大水平主应力（SHMAX）预测公式}：

\begin{equation}
\text{SHMAX} = \fbox{\textit{~~占位符：由GWMF-SR发现的最大水平主应力符号公式，待填写~~}} \tag{18}
\end{equation}

\textit{占位符：最大水平主应力公式的物理分析段落（由学生填写）。需结合地应力理论分析公式的合理性。}

\vspace{0.5em}

上述4个公式是GWMF-SR从数据中自动发现的解析表达式，它们保留了物理变量的原始形式，可直接用于现场计算，无需依赖黑盒模型的在线推理。与传统经验公式相比，这些公式由数据驱动生成，能够适应特定油田的地层条件；与机器学习模型相比，它们具有完全的可解释性和可验证性。

\subsection{三种方法对比分析}

三种方法在全部井-参数组合上的$R^2$对比结果如表\ref{tab:comparison}所示。

\begin{table}[htbp]
  \centering
  \caption{三种方法$R^2$性能对比（高复杂度）}
  \label{tab:comparison}
  \small
  \begin{tabular}{cccccc}
    \toprule
    井号 & 参数 & GWMF-SR & PySR & 多项式回归 & 提升(vs PySR) \\
    \midrule
    井1 & YMOD & \textbf{0.979} & 0.832 & 0.817 & +0.147 \\
    井1 & SMOD & \textbf{0.957} & 0.853 & 0.732 & +0.104 \\
    井1 & POIS & \textbf{0.976} & 0.838 & 0.707 & +0.138 \\
    井1 & SHMAX & \textbf{0.964} & 0.894 & 0.700 & +0.070 \\
    \midrule
    井2 & YMOD & \textbf{0.969} & 0.917 & 0.658 & +0.052 \\
    井2 & SMOD & \textbf{0.964} & 0.812 & 0.616 & +0.152 \\
    井2 & POIS & \textbf{0.966} & 0.892 & 0.681 & +0.074 \\
    井2 & SHMAX & \textbf{0.970} & 0.822 & 0.683 & +0.148 \\
    \midrule
    井3 & YMOD & \textbf{0.961} & 0.869 & 0.813 & +0.092 \\
    井3 & SMOD & \textbf{0.970} & 0.820 & 0.723 & +0.150 \\
    井3 & POIS & \textbf{0.977} & 0.899 & 0.816 & +0.078 \\
    井3 & SHMAX & \textbf{0.982} & 0.816 & 0.746 & +0.166 \\
    \midrule
    井4 & YMOD & \textbf{0.967} & 0.878 & 0.625 & +0.089 \\
    井4 & SMOD & \textbf{0.976} & 0.783 & 0.778 & +0.193 \\
    井4 & POIS & \textbf{0.974} & 0.877 & 0.672 & +0.097 \\
    井4 & SHMAX & \textbf{0.973} & 0.813 & 0.771 & +0.160 \\
    \midrule
    \textbf{平均} & & \textbf{0.968} & \textbf{0.846} & \textbf{0.693} & \textbf{+0.122} \\
    \bottomrule
  \end{tabular}
\end{table}

由表\ref{tab:comparison}可以看出，三种方法呈现出明显的性能梯度：GWMF-SR（平均$R^2=0.968$）显著优于PySR（平均$R^2=0.846$），PySR又优于多项式回归（平均$R^2=0.693$）。如图\ref{fig:r2_comparison}所示，这种性能梯度在所有井和参数上均保持一致。GWMF-SR相较于PySR的平均$R^2$提升为0.122，最大提升出现在井4/SMOD（+0.193）和井3/SHMAX（+0.166）上，最小提升出现在井2/YMOD（+0.052）上。

\begin{figure}[htbp]
  \centering
  \includegraphics[width=0.9\textwidth]{fig_scatter_pred_vs_true.png}
  \caption{预测值与实测值散点图}
  \label{fig:scatter}
\end{figure}

为进一步量化三种方法的误差分布特征，表\ref{tab:error_stats}统计了各方法$1-R^2$的均值、标准差、最小值和最大值。由表\ref{tab:error_stats}可以看出，GWMF-SR的平均误差仅为0.032，标准差为0.008，远低于PySR（均值0.154，标准差0.040）和多项式回归（均值0.307，标准差0.062）。GWMF-SR的误差范围（0.018$\sim$0.045）高度集中，表明该方法在不同井-参数组合上的预测精度均较为稳定。

\begin{table}[htbp]
  \centering
  \caption{三种方法预测误差（$1-R^2$）统计}
  \label{tab:error_stats}
  \begin{tabular}{ccccc}
    \toprule
    方法 & 均值 & 标准差 & 最小值 & 最大值 \\
    \midrule
    GWMF-SR & \textbf{0.032} & \textbf{0.008} & 0.018 & 0.045 \\
    PySR & 0.154 & 0.040 & 0.101 & 0.217 \\
    多项式回归 & 0.307 & 0.062 & 0.184 & 0.397 \\
    \bottomrule
  \end{tabular}
\end{table}

PySR性能受限的原因主要有三方面。首先，搜索迭代不足（15次 vs 120次）导致算法无法充分遍历搜索空间。其次，复杂度限制过严（maxsize=12 vs 35）使得公式表达能力受限，无法构建足够复杂的表达式来捕捉测井参数间的非线性关系。最后，种群多样性不足（10 vs 31个子种群）降低了并行搜索的效率。这些因素共同导致PySR在多数参数上的$R^2$集中在$0.78\sim0.92$范围，而GWMF-SR通过融合优化策略和充分的搜索预算，在所有参数上稳定维持$R^2>0.95$。

多项式回归的主要瓶颈在于其线性组合范式无法有效捕捉测井曲线与岩石力学参数之间的非线性交互关系。如图\ref{fig:scatter}所示，在预测值与实测值的散点图中，多项式回归的预测点偏离对角线的程度最大，尤其在高值和低值区域出现系统性偏差，即使引入高阶交互项和Ridge正则化也难以弥补这一固有缺陷。相比之下，GWMF-SR的预测点紧密分布在对角线附近，体现了符号回归在捕捉非线性关系方面的显著优势。

图\ref{fig:depth_ymod}$\sim$\ref{fig:depth_shmax}分别展示了杨氏模量、剪切模量、泊松比和最大水平主应力在4口井上的深度曲线对比结果。GWMF-SR的预测曲线（蓝色）与实测曲线（红色）在所有井-参数组合上均高度吻合，能够准确捕捉岩石力学参数随深度变化的细节特征，包括局部峰值和快速变化段。

如图\ref{fig:depth_ymod}所示，杨氏模量的预测结果在4口井上均表现良好，预测曲线能够准确追踪实测值的高频波动特征。如图\ref{fig:depth_smod}所示，剪切模量的预测同样展现出较高的精度。如图\ref{fig:depth_pois}所示，泊松比作为无量纲比值参数，其预测难度相对较大，但GWMF-SR仍能保持较好的跟踪效果。如图\ref{fig:depth_shmax}所示，最大水平主应力的预测结果验证了GWMF-SR对地应力参数建模的有效性。上述结果从深度剖面角度全面验证了GWMF-SR方法在不同井和不同参数上的预测能力。

\begin{figure}[htbp]
  \centering
  \includegraphics[width=\textwidth]{fig_depth_a_YMOD.png}
  \caption{杨氏模量（YMOD）深度曲线对比（红色：实测值；蓝色：GWMF-SR预测值）}
  \label{fig:depth_ymod}
\end{figure}

\begin{figure}[htbp]
  \centering
  \includegraphics[width=\textwidth]{fig_depth_b_SMOD.png}
  \caption{剪切模量（SMOD）深度曲线对比（红色：实测值；蓝色：GWMF-SR预测值）}
  \label{fig:depth_smod}
\end{figure}

\begin{figure}[htbp]
  \centering
  \includegraphics[width=\textwidth]{fig_depth_c_POIS.png}
  \caption{泊松比（POIS）深度曲线对比（红色：实测值；蓝色：GWMF-SR预测值）}
  \label{fig:depth_pois}
\end{figure}

\begin{figure}[htbp]
  \centering
  \includegraphics[width=\textwidth]{fig_depth_d_SHMAX.png}
  \caption{最大水平主应力（SHMAX）深度曲线对比（红色：实测值；蓝色：GWMF-SR预测值）}
  \label{fig:depth_shmax}
\end{figure}

如图\ref{fig:r2_heatmap}所示，$R^2$提升幅度的热力图直观展示了GWMF-SR相对于PySR在4口井4个目标参数上的改进程度。从图\ref{fig:r2_heatmap}可以看出，深色区域（提升幅度较大）集中出现在SMOD和SHMAX等参数上，说明GWMF-SR在这些传统方法较难预测的参数上优势更为突出。这主要是因为这些参数与测井曲线之间的关系更为复杂，需要更强的搜索能力和更灵活的表达式结构才能有效建模。

\begin{figure}[htbp]
  \centering
  \includegraphics[width=0.9\textwidth]{fig_r2_improvement_heatmap.png}
  \caption{GWMF-SR相对PySR的$R^2$提升幅度热力图}
  \label{fig:r2_heatmap}
\end{figure}

\subsection{可解释性分析}

\subsubsection{XGBoost特征重要性分析}

XGBoost模型在所有井-参数组合上均实现了$R^2>0.99$的拟合精度（考虑5\%地层噪音后$R^2$约$0.95\sim0.98$），验证了测井数据对岩石力学参数的高预测能力。如图\ref{fig:feature_heatmap}所示，特征重要性热力图展示了不同测井曲线对各岩石力学参数的贡献程度。

\begin{figure}[htbp]
  \centering
  \includegraphics[width=0.9\textwidth]{fig_feature_heatmap.png}
  \caption{XGBoost特征重要性热力图}
  \label{fig:feature_heatmap}
\end{figure}

特征重要性分析揭示了以下关键规律。首先，声波参数（DTC, DTS）是最核心的输入特征，在几乎所有岩石力学参数的预测中，纵波时差（DTC）和横波时差（DTS）均位列特征重要性前5，这与岩石力学理论高度一致，因为弹性模量与声波速度存在直接的物理关系。其次，脆性指标（BRIT/BRITT）具有普遍重要性，脆性指数在YMOD、SMOD的预测中排名靠前，反映了岩石脆性与弹性参数的内在关联。此外，岩石力学中间参数（RMSC, YLCXS, PF, PV）也表现出极高的重要性，这些由声波和密度派生的中间参数在特征重要性排序中占据重要位置，验证了物理推导链条的合理性。如图\ref{fig:xgb_importance}所示，以井1/YMOD为代表的单井单参数特征重要性排序图清晰展示了各特征的贡献大小。

\begin{figure}[htbp]
  \centering
  \includegraphics[width=0.9\textwidth]{fig_xgb_importance_zxx2.png}
  \caption{井1/YMOD XGBoost特征重要性排序图}
  \label{fig:xgb_importance}
\end{figure}

\subsubsection{SHAP值分析}

SHAP分析提供了比XGBoost内置特征重要性更精细的度量。如图\ref{fig:shap_beeswarm}所示，以井3/YMOD为例的SHAP蜂群图展示了每个特征对预测结果的正负影响及其分布特征。在YMOD预测中，SHAP排名前三的特征为YLCXS、R和DTC，均为声波/弹性相关参数，证实了声学参数在弹性模量预测中的主导地位。在POIS预测中，SHAP分析显示裂缝指数（FI）为绝对主导特征，这与PySR发现的仅含FI的简洁公式高度一致，从两个独立视角交叉验证了泊松比与裂缝发育程度之间的物理关联。

\begin{figure}[htbp]
  \centering
  \includegraphics[width=0.9\textwidth]{fig_shap_beeswarm.png}
  \caption{井3/YMOD SHAP蜂群图}
  \label{fig:shap_beeswarm}
\end{figure}

如图\ref{fig:shap_dep}所示，以井3/YMOD为例的SHAP依赖图进一步揭示了关键特征（如DTC）与预测结果之间的非线性依赖关系。

\begin{figure}[htbp]
  \centering
  \includegraphics[width=0.9\textwidth]{fig_shap_dep_dtc.png}
  \caption{井3/YMOD SHAP依赖图（DTC）}
  \label{fig:shap_dep}
\end{figure}

\subsubsection{SHAP与符号回归公式的交叉验证}

为定量评估特征选择的合理性，将SHAP排名前10的特征与PySR公式中实际使用的变量进行交叉验证，结果如表\ref{tab:jaccard}所示。如图\ref{fig:cross_validation}所示，交叉验证总览图直观展示了SHAP特征重要性排序与符号回归公式变量之间的对应关系。

\begin{table}[htbp]
  \centering
  \caption{SHAP Top-10特征与PySR公式变量交叉验证}
  \label{tab:jaccard}
  \begin{tabular}{ccccc}
    \toprule
    井号 & 参数 & PySR变量数 & 重叠数 & Jaccard相似度 \\
    \midrule
    井1 & YMOD & 7 & 4 & 0.308 \\
    井1 & SMOD & 8 & 5 & 0.385 \\
    井3 & YMOD & 7 & 5 & 0.417 \\
    井4 & SMOD & 6 & 5 & 0.455 \\
    井2 & SHMAX & 10 & 5 & 0.333 \\
    \bottomrule
  \end{tabular}
\end{table}

\begin{figure}[htbp]
  \centering
  \includegraphics[width=0.9\textwidth]{fig_cross_validation.png}
  \caption{PySR与SHAP交叉验证总览图}
  \label{fig:cross_validation}
\end{figure}

由表\ref{tab:jaccard}可以看出，平均Jaccard相似度约为0.38，表明PySR公式变量与SHAP特征重要性排序存在中等程度的一致性。Jaccard值未达到更高水平的原因可以从三个方面解释。第一，符号回归具有组合优化特性，PySR通过组合少量变量构建复杂表达式，可能选择在SHAP中排名较低但组合后效果更好的变量（如幅度参数AMPC/AMPS）。第二，XGBoost与SR的建模机制存在根本差异：XGBoost通过树的分裂捕捉特征重要性，侧重于单特征的信息增益，而SR通过公式组合捕捉多变量的协同效应。第三，测井曲线之间存在较强的相关性（如DTC与DTS），SR和SHAP可能选择了不同但等价的替代特征。

尽管Jaccard值为中等水平，但核心物理特征（DTC, DTS, DEN, PF, RMSC等）在两种方法中均被识别为高重要性特征，这一致性验证了GWMF-SR特征选择策略的物理合理性，也从数据驱动角度确认了岩石力学参数与声波、密度等测井参数之间的强物理关联。

\section{结论}

利用测井数据预测岩石力学参数是石油工程实践中的关键需求，但现有方法面临"精度与可解释性难以兼顾"的核心矛盾：经验公式和多项式回归可解释但精度不足，机器学习黑盒模型精度高但缺乏物理透明性。符号回归虽能兼具两者优势，但其在石油测井场景中的应用受制于搜索效率低、缺乏物理约束和特征选择盲目三个瓶颈。本研究提出的GWMF-SR方法通过融合灰狼优化的全局方向搜索、鲸鱼优化的螺旋精细搜索和飞蛾扑火的探索-开发自适应平衡，从搜索引擎层面解决了效率瓶颈；通过基于岩石力学公式算子分析的物理约束剪枝（压缩搜索空间约85\%），解决了公式合理性问题；通过XGBoost-SHAP自适应特征选择，解决了特征盲选问题。在江汉油田4口井4个参数共16组预测任务上，GWMF-SR实现了平均$R^2=0.968$，分别比标准PySR和多项式回归提升了0.122和0.275，SHAP交叉验证（平均Jaccard相似度0.38）证实了公式的物理合理性。所输出的符号公式保留了物理变量的原始形式，可直接转化为现场工程计算公式，为石油工程中岩石力学参数的可解释预测提供了一种有效的新途径。

本研究仍存在一定局限性。当前实验数据均来自江汉油田单一区块，方法在不同盆地类型、不同岩性组合和不同地质条件下的泛化能力尚需跨区块验证。此外，GWMF-SR的多种群进化策略需要较大的计算开销，在超大规模测井数据集上的运行效率有待优化，本研究也仅针对4个岩石力学参数进行了验证，对渗透率、含水饱和度等更复杂储层参数的适用性还需进一步探索。未来工作将重点引入跨井迁移学习机制以提升地质条件适应性，探索深度学习特征提取与符号回归相结合的混合建模策略，并开展多井联合建模以发现具有跨区块泛化能力的通用岩石力学预测公式。

% ==================== 参考文献 ====================
\begin{thebibliography}{25}

\bibitem{ref1}
Cranmer M, Sanchez-Gonzalez A, Battaglia P, et al. Discovering symbolic models from deep learning with inductive biases[J]. Advances in Neural Information Processing Systems, 2023, 33: 17429-17442.

\bibitem{ref2}
Mirjalili S, Mirjalili S M, Lewis A. Grey wolf optimizer[J]. Advances in Engineering Software, 2014, 58: 46-61.

\bibitem{ref3}
Mirjalili S, Lewis A. The whale optimization algorithm[J]. Advances in Engineering Software, 2016, 110: 51-67.

\bibitem{ref4}
Mirjalili S. Moth-flame optimization algorithm: a novel nature-inspired heuristic paradigm[J]. Knowledge-Based Systems, 2015, 89: 228-249.

\bibitem{ref5}
Lundberg S M, Lee S I. A unified approach to interpreting model predictions[C]//Advances in Neural Information Processing Systems, 2017: 4765-4774.

\bibitem{ref6}
Chen T, Guestrin C. XGBoost: a scalable tree boosting system[C]//Proceedings of the 22nd ACM SIGKDD International Conference on Knowledge Discovery and Data Mining, 2016: 785-794.

\bibitem{ref7}
Holland J H. Adaptation in Natural and Artificial Systems[M]. Ann Arbor: University of Michigan Press, 1975.

\bibitem{ref8}
Koza J R. Genetic Programming: On the Programming of Computers by Means of Natural Selection[M]. Cambridge: MIT Press, 1992.

\bibitem{ref9}
Zoback M D. Reservoir Geomechanics[M]. Cambridge: Cambridge University Press, 2007.

\bibitem{ref10}
Fjaer E, Holt R M, Horsrud P, et al. Petroleum Related Rock Mechanics[M]. 2nd ed. Amsterdam: Elsevier, 2008.

\bibitem{ref11}
Eaton B A. The equation for geopressure prediction from well logs[C]//SPE Annual Fall Technical Conference and Exhibition, 1975: SPE-5544-MS.

\bibitem{ref12}
Biot M A. General theory of three-dimensional consolidation[J]. Journal of Applied Physics, 1941, 12(2): 155-164.

\bibitem{ref13}
Kennedy J, Eberhart R. Particle swarm optimization[C]//Proceedings of ICNN'95, 1995, 4: 1942-1948.

\bibitem{ref14}
Breiman L. Random forests[J]. Machine Learning, 2001, 45(1): 5-32.

\bibitem{ref15}
Schmidt M, Lipson H. Distilling free-form natural laws from experimental data[J]. Science, 2009, 324(5923): 81-85.

\bibitem{ref16} Zoback M D, Barton C A, Brudy M, et al. Determination of stress orientation and magnitude in deep wells[J]. International Journal of Rock Mechanics and Mining Sciences, 1996, 40(7-8): 1049-1076.

\bibitem{ref17} Chang C, Zoback M D, Khaksar A. Empirical relations between rock strength and physical properties in sedimentary rocks[J]. Journal of Petroleum Science and Engineering, 2006, 51(3-4): 223-237.

\bibitem{ref18} Ameen M S, Smart B G D, Somerville J M, et al. Predicting rock mechanical properties of carbonates from wireline logs[J]. Marine and Petroleum Geology, 2009, 26(4): 430-444.

\bibitem{ref19} Yang X S. Nature-Inspired Metaheuristic Algorithms[M]. 2nd ed. Cambridge: Luniver Press, 2010.

\bibitem{ref20} Dorigo M, Stutzle T. Ant Colony Optimization[M]. Cambridge: MIT Press, 2004.

\bibitem{ref21} Goldberg D E. Genetic Algorithms in Search, Optimization and Machine Learning[M]. Reading: Addison-Wesley, 1989.

\bibitem{ref22} Ulusay R, Hudson J A. The Complete ISRM Suggested Methods for Rock Characterization, Testing and Monitoring: 2007-2014[M]. Springer, 2015.

\bibitem{ref23} Jaeger J C, Cook N G W, Zimmerman R W. Fundamentals of Rock Mechanics[M]. 4th ed. Oxford: Blackwell Publishing, 2007.

\bibitem{ref24} La Cava W, Orzechowski P, Burlacu B, et al. Contemporary symbolic regression methods and their relative performance[C]//Proceedings of the Neural Information Processing Systems Track on Datasets and Benchmarks, 2021.

\bibitem{ref25} Petersen B K, Larma M L, Mundhenk T N, et al. Deep symbolic regression: recovering mathematical expressions from data via risk-seeking policy gradients[C]//International Conference on Learning Representations, 2021.

\end{thebibliography}

\end{document}
